\section{An endpoint certificate uniform up to critical corners}
\label{sec:v103-certificate}
Write $\theta=\theta_0+\sigma x$, $0\le x\le L_\theta$, for
one angular side, where $\sigma\in\{-1,1\}$. Derivatives below
are with respect to $x$. A certificate is imposed on the actual
regular pairing chamber, not on a symbolic continuation of the word.

\begin{theorem}[An expanding endpoint-distance certificate]
\label{thm:v103-certificate}
Choose either $e=a$ or $e=b$. Suppose the chamber continues over
this angular interval, $c^\varepsilon>0$, and for constants
$\nu>0$, $A,K\ge0$ one has
\begin{equation}\label{eq:v103-certificate-assumptions}
 |\dot e|\ge\nu,\qquad
 0\le\dot\ell\le K|\dot e|,\qquad
 |(\log c^\varepsilon)^{\boldsymbol\cdot}|\le A,
 \qquad |(\log|\dot e|)^{\boldsymbol\cdot}|\le A
\end{equation}
almost everywhere. Let $0<\ell(\theta_0)\le\delta$ and assume
\begin{equation}\label{eq:v103-certificate-scale}
 r\le\nu L_\theta,\qquad
 2Ar/\nu\le\log2,\qquad \delta+Kr<\Delta.
\end{equation}
Then the clock for the endpoint $e$ satisfies
\[
                  k_{r,\Delta}^{e}(\theta_0,s)\ge\frac12
                  \quad\text{for every }0<s<1.
\]
Here $k^a=k^-$ and $k^b=k^+$. The conclusion uses no lower bound
on $\ell(\theta_0)$ or on $s,1-s$, and no bound on
$\partial_\theta\log\ell$. The sign of $\dot e$ may be either
positive or negative.
\end{theorem}
\begin{proof}
The nonzero continuous derivative has a constant sign $\tau$.
For every $0\le h\le r$, the equation
$e(\theta_h)-e(\theta_0)=\tau h$ has a unique solution on the
chosen side, with angular displacement at most $h/\nu$.
Choose the gap to that endpoint at the starting state and keep it
fixed. Initially it lies strictly between zero and $\ell(\theta_0)$.
The residence is nondecreasing on the chosen side and
\[
 \ell(\theta_0)\le\ell(\theta_h)
 \le\ell(\theta_0)+K|e(\theta_h)-e(\theta_0)|
 \le\delta+Kr<\Delta.
\]
The fixed gap therefore stays strictly inside the same segment.
There is no loss at its other endpoint, however close the starting
relative coordinate was to that endpoint. The regularity assumption
and positive weight keep the rest of the comparison domain intact.
This curve is exactly the partial flow at roof time $\tau h$.
By \eqref{eq:v103-gap-jacobian} and the two logarithmic bounds,
\[
 J_{\tau h}^e\ge\exp(-2Ah/\nu)\ge\tfrac12.
\]
One full one-sided time interval of length $r$ now contributes at
least $1/2$ to \eqref{eq:v103-gap-clock}. Endpoint values of the
angular interval are handled by smaller times and monotone convergence.
\end{proof}

\begin{corollary}[Critical entrances and opposite endpoint slopes]
\label{cor:v103-corners}
On a regular backward-critical chamber with $a$ constant, choose
an angular side on which $b$ increases. Then $\dot\ell=\dot b$
and Theorem~\ref{thm:v103-certificate} applies with $e=b$, $K=1$,
provided its speed, weight and reach hypotheses hold. The conclusion
is uniform as $s\downarrow0$, although the relative-clock speed is
$v=sb'$.

On a side with $\dot a\le0\le\dot b$, if one endpoint derivative
has magnitude at least the other throughout the comparison interval,
use that endpoint as $e$. Then
$0\le\dot\ell\le2|\dot e|$. With $\Delta=4\delta$ and $r=\delta$
the residence constraint holds with $K=2$. Subject to the same
remaining hypotheses, every $s$ is admitted, including a value at
which $(1-s)a'+sb'=0$.
\end{corollary}
\begin{proof}
For the first assertion differentiate $b-a$. For the second,
$\dot\ell=\dot b+|\dot a|$ is at most twice the larger endpoint
speed. The certificate and $\delta+2\delta<4\delta$ finish the
proof. No sign is assigned to the positive incoming critical trace,
which is explicitly given in Proposition~\ref{prop:v103-critical-trace}.
\end{proof}

For any Borel physical subfamily on which this certificate holds
with $\Delta=4\delta,r=\delta$, Theorem~\ref{thm:v103-gap-height}
can be used with threshold $1/2$. Its complete positive trace has
ordered height
\begin{equation}\label{eq:v103-certificate-height}
 C\{\varepsilon^{1/16}+(1+\varepsilon^{-1})\delta\}
                                      =O(B^{-1/192}).
\end{equation}
The estimate is independent of how many physical chambers satisfy
the tests. This is a height assertion for the certified physical
source, not just a sheetwise derivative bound.

\subsection{Two checks of the newly recovered geometry}
The following elementary chamber models test the comparison formulas.
They are not asserted to be additional Lorentz realizations. Take
$c^\varepsilon=1$, $r=\delta$, and a collar $0<\ell<4\delta$.
If $a=a_*$ and $b=a_*+\theta$, then the outgoing-gap flow is
\[
 \theta_h=\theta+h,\qquad u_h=u+h,\qquad J_h^+=1.
\]
For $\ell(\theta)\le\delta$ its whole forward interval $0<h<r$
exists, at every starting relative position. In contrast, at fixed
$s$ the relative flow has total roof reach at most $4s\delta$,
so its capped clock is at most $4s$. Thus for any fixed positive
admission threshold there is a positive-measure critical-corner
strip missed by the relative clock and recovered by the gap clock.

For $a=-\theta$, $b=\theta$, $\theta>0$, the outgoing-gap flow
again has $\theta_h=\theta+h$, $u_h=u+h$ and $J_h^+=1$.
Its residence increases by $2h$, so all initial $\ell\le\delta$
survive the full forward comparison. The relative speed is $2s-1$;
its total roof reach within the collar is at most
$2\delta|2s-1|$. It is zero at the fold and arbitrarily small on
a positive-measure neighborhood, while the gap clock remains at
least one. The tests also show why neither the $s$ factor nor the
relative residence derivative belongs in the new Jacobian.

\subsection{The remaining full-source statement}
The recovered height bounds do not verify the certificate on every
physical pair. In particular simultaneous stationary endpoint roofs,
flat-interval pulses from Lemma~\ref{lem:v103-flat}, insufficient
one-sided comparison reach and rapid variation of the positive
minimum may remain in $\mathsf U_B^{\rm rem}$. The proof does not
set a nonregular or backward-critical trace to zero. Establishing
$\mathcal H_M(\mathsf U_B^{\rm rem})\to0$ on the full Lorentz
family remains the required upper pointwise step. The original
pointwise objective and all scalar-first arithmetic restrictions
are unchanged.
